\section{Orbital angular-momentum density: the axial component}
\label{sec:canonical-oam-density-and_flux}

We consider the $Oz$ component, with $(x,y,z)=(x_1,x_2,x_3)$ and
${\bf r}={\bf x}$.  Tildes denote Fourier amplitudes at $(\omega,{\bf x})$,
where $\omega>0$.  In the radiation gauge,
$\tilde{\bf A}=\tilde{\bf E}/(i\omega)$.
We use the electric-field canonical expression.  In spectral expressions,
$\mathcal L_z$ denotes the time-integrated density.  The one-sided bilinear
spectral factor is $4\pi$ for the Fourier normalization used in this document.

The canonical orbital angular momentum density in the time domain is:
\begin{align}
  {\boldsymbol{\mathcal L}} = \epsilon_0 \sum_{i=x,y,z} E_i\, ({\bf r}\times{\boldsymbol{\nabla}})\, A_i
\end{align}

Its $z$-component involves the transverse differential operator $\hat{L}_z = ({\bf r}\times{\boldsymbol{\nabla}})_z = (x\partial_y - y\partial_x)$:
\begin{align}
  {\cal L}_z = \epsilon_0 \sum_{i=x,y,z} E_i\, (x\partial_y - y\partial_x) A_i
\end{align}

With the Fourier convention of Eq.~\eqref{eq:Faraday-Fourier-convention-observables}, the one-sided spectrum of the time-integrated density is:
\begin{align}
  \frac{d {\cal L}_z}{d\omega} = 4\pi\epsilon_0 \sum_{i=x,y,z} \operatorname{Re}\left[\tilde E_i\, \Big(\hat{L}_z \tilde A_i\Big)^*\right]
\end{align}

Using $\tilde A_i^* = \frac{i}{\omega}\tilde E_i^*$ and noting that coordinates and spatial derivative operators are real:
\begin{align}
  \Big(\hat{L}_z \tilde A_i\Big)^* = \frac{i}{\omega}\hat{L}_z \tilde E_i^*
\end{align}
\begin{align}
  \operatorname{Re}\left[\tilde E_i \left(\frac{i}{\omega}\hat{L}_z \tilde E_i^*\right)\right] = -\frac{1}{\omega}\operatorname{Im}\left[\tilde E_i \hat{L}_z \tilde E_i^*\right] = \frac{1}{\omega}\operatorname{Im}\left[\tilde E_i^*\, \hat{L}_z \tilde E_i\right]
\end{align}

\subsection{Rectangular screen: Cartesian grid}
On a Cartesian grid $(x, y)$, the operator is $\hat{L}_z = x \partial_y - y \partial_x$:
\begin{importantbox}
  \begin{align}
    \frac{d {\cal L}_z}{d\omega} = \frac{4\pi\epsilon_0}{\omega}\sum_{i=x,y,z}\operatorname{Im}\left[\tilde E_i^* \left(x \frac{\partial \tilde E_i}{\partial y} - y \frac{\partial \tilde E_i}{\partial x}\right)\right] \label{eq:canonical-oam-density-cartesian}
  \end{align}
\end{importantbox}

\subsection{Circular screen: polar grid}
On a polar grid $(\rho, \phi)$ where $x = \rho\cos\phi$, $y = \rho\sin\phi$, we use the identity $x\partial_y - y\partial_x = \partial_\phi$:
\begin{importantbox}
  \begin{align}
    \frac{d {\cal L}_z}{d\omega} = \frac{4\pi\epsilon_0}{\omega}\sum_{i=x,y,z}\operatorname{Im}\left[\tilde E_i^* \frac{\partial \tilde E_i}{\partial \phi}\right] \label{eq:canonical-oam-density-polar}
  \end{align}
\end{importantbox}

The polar derivative acts on Cartesian field components at fixed $\rho$ and $z$.

\section{Orbital angular-momentum flux and continuity equation}
\label{sec:canonical-oam-flux}

The local conservation law for orbital angular momentum in vacuum is
\begin{importantbox}
  \begin{align}
    \partial_t\mathcal L_i+\partial_j\Lambda_{ij}=0.
    \label{eq:canonical-oam-continuity}
  \end{align}
\end{importantbox}
Here $\mathcal L_i=\epsilon_0\sum_m E_m
({\bf r}\times\boldsymbol{\nabla})_i A_m$ is the time-domain canonical
orbital angular-momentum density defined above.  Repeated spatial indices
are summed, and $\Lambda_{ij}$ denotes the flux of its $i$th component in
the $j$th direction.  In SI units, the flux tensor is
\begin{align}
  \Lambda_{ij}
  =\epsilon_0c^2\left\{
    \varepsilon_{ikl}r_k\left[
      \varepsilon_{jmn}B_n(\partial_l A_m)
      +\frac12\delta_{lj}\left(\frac{E^2}{c^2}-B^2\right)
    \right]+B_j A_i
  \right\}.
  \label{eq:canonical-oam-flux-tensor}
\end{align}

\subsection{Reduction of the axial component}

For $i=j=z=3$, the trace term vanishes because
$\varepsilon_{3k3}=0$.  Using the transverse rotation operator
$\hat L_z=x\partial_y-y\partial_x$, the time-domain axial flux becomes
\begin{highlightbox}
  \begin{align}
    \Lambda_{zz}=\epsilon_0c^2\left[
      B_y\hat L_z A_x-B_x\hat L_z A_y+B_z A_z
    \right].
    \label{eq:canonical-oam-axial-flux}
  \end{align}
\end{highlightbox}

\subsection{Spectral density of the axial flux}

As for the density, $d\Lambda_{zz}/d\omega$ denotes the one-sided
spectrum of the time-integrated flux.  Applying the bilinear spectral
rule with the document's Fourier normalization gives
\begin{align}
  \frac{d\Lambda_{zz}}{d\omega}
  =4\pi\epsilon_0c^2\operatorname{Re}\left[
    \tilde B_y(\hat L_z\tilde A_x)^*
    -\tilde B_x(\hat L_z\tilde A_y)^*
    +\tilde B_z\tilde A_z^*
  \right].
  \label{eq:canonical-oam-flux-bilinear}
\end{align}
Substituting $\tilde A_k^*=i\tilde E_k^*/\omega$ and using the reality of
$\hat L_z$ yields
\begin{align}
  \frac{d\Lambda_{zz}}{d\omega}
  =\frac{4\pi\epsilon_0c^2}{\omega}\operatorname{Im}\left[
    \tilde B_y^*\hat L_z\tilde E_x
    -\tilde B_x^*\hat L_z\tilde E_y
    +\tilde B_z^*\tilde E_z
  \right],\qquad \omega>0.
  \label{eq:canonical-oam-flux-spectrum}
\end{align}

\subsubsection{Rectangular screen: Cartesian grid}

On the Cartesian grid, $\hat L_z=x\partial_y-y\partial_x$, so
\begin{importantbox}
  \begin{align}
    \frac{d\Lambda_{zz}}{d\omega}
    =\frac{4\pi\epsilon_0c^2}{\omega}\operatorname{Im}\Bigg[&
      \tilde B_y^*\left(
        x\frac{\partial\tilde E_x}{\partial y}
      -y\frac{\partial\tilde E_x}{\partial x}\right)
      \nonumber\\
      &-\tilde B_x^*\left(
        x\frac{\partial\tilde E_y}{\partial y}
      -y\frac{\partial\tilde E_y}{\partial x}\right)
    +\tilde B_z^*\tilde E_z\Bigg].
    \label{eq:canonical-oam-flux-cartesian}
  \end{align}
\end{importantbox}

\subsubsection{Circular screen: polar grid}

On the polar grid, $\hat L_z=\partial_\phi$, and hence
\begin{importantbox}
  \begin{align}
    \frac{d\Lambda_{zz}}{d\omega}
    =\frac{4\pi\epsilon_0c^2}{\omega}\operatorname{Im}\left[
      \tilde B_y^*\frac{\partial\tilde E_x}{\partial\phi}
      -\tilde B_x^*\frac{\partial\tilde E_y}{\partial\phi}
      +\tilde B_z^*\tilde E_z
    \right].
    \label{eq:canonical-oam-flux-polar}
  \end{align}
\end{importantbox}
Both expressions apply for $\omega>0$.  All Fourier amplitudes are evaluated
at $(\omega,{\bf x})$, and the polar derivative acts on Cartesian field
components at fixed $\rho$ and $z$.

\subsection{Main results in terms of the Faraday tensor}

Using Eq.~\eqref{eq:Fourier-fields-from-Faraday}, the Fourier amplitudes
are related to the Faraday tensor by
\begin{align}
  \tilde E_i=cF^{i0},\qquad
  \tilde B_x=F^{32},\qquad
  \tilde B_y=F^{13},\qquad
  \tilde B_z=F^{21}.
\end{align}
Here $i=1,2,3$ corresponds to $x,y,z$, and every tensor component below
is evaluated at $(\omega,{\bf x})$.

\begin{highlightbox}[title={Orbital angular-momentum density and axial flux from the Faraday tensor}]
  The one-sided spectra of the time-integrated density and flux are
  \begin{align}
    \frac{d\mathcal L_z}{d\omega}
    &=\frac{4\pi\epsilon_0c^2}{\omega}
    \operatorname{Im}\left[
      \sum_{i=1}^{3}(F^{i0})^*\hat L_z F^{i0}
    \right],
    \label{eq:canonical-oam-density-Faraday}\\
    \frac{d\Lambda_{zz}}{d\omega}
    &=\frac{4\pi\epsilon_0c^3}{\omega}
    \operatorname{Im}\left[
      (F^{13})^*\hat L_z F^{10}
      -(F^{32})^*\hat L_z F^{20}
      +(F^{21})^*F^{30}
    \right].
    \label{eq:canonical-oam-flux-Faraday}
  \end{align}
  Both results hold for $\omega>0$, with
  \begin{align}
    \hat L_z=x\frac{\partial}{\partial y}
    -y\frac{\partial}{\partial x}
    =\frac{\partial}{\partial\phi}.
  \end{align}
  The Cartesian form applies on a rectangular grid; the azimuthal form
  applies on a polar grid at fixed $\rho$ and $z$.  The derivatives act
  on the Cartesian tensor components immediately to their right.
\end{highlightbox}
